\section{Empirical Application}\label{sec:6}

I apply the estimators to Indonesia's INPRES school construction program, a canonical setting with imperfect compliance. School construction varied across districts and cohorts, but exposure to the program increased schooling without deterministically fixing completed education. \citet{duflo2001schooling} study this design with 2SLS, and \citet{fuzzy2018} recast it as a discrete DID-IV setting. This application is a useful test case for the DML estimators because the first-stage coefficient is small, the covariates are rich, and the propensity scores approach the boundary. I compare DML-Wald and DML-TC under symmetric trimming and the $\hat{\alpha}$-rule.

\subsection{Setting, data, and design}\label{sec:6_setup}

The INPRES program built more than 61{,}000 primary schools between 1973 and 1978, creating variation in school access across districts and cohorts. I use the 1995 Indonesian Intercensal Survey (SUPAS) microdata from \citet{duflo2001schooling}, pooling men and women with non-missing log hourly wage born in 1957--1962 (old cohort, $T=0$) or 1968--1972 (young cohort, $T=1$).

I adopt a binary district classifier $G^*$, which flags districts whose schooling distribution shifted upward between cohorts via a signed chi-squared test on the district's cohort-by-education table (Appendix~\ref{sec:app_classifier} provides details). \citet{fuzzy2018} also retain those districts with downward shifts by pooling $G^* \in \{-1, 0, +1\}$, and I exclude these 20 districts from the analysis. The final sample has $22{,}414$ observations in $201$ districts.

I compare two binary treatment thresholds, $D=\mathbf{1}\{\text{yeduc} \geq 6\}$ (primary completion) and $D=\mathbf{1}\{\text{yeduc} \geq 12\}$ (high-school completion), instrumented by the group-cohort interaction $G \times T$. The $p = 45$ covariate dictionary contains only individual-level regressors (sex, birth year interacted with cohort, Java, urban, and pairwise interactions). District-level covariates that mechanically predict $G^*$ are excluded because they would push $\hat e(X)$ to one and leave holdout cells empty. Appendix~\ref{sec:app_empirical} gives the full construction.


\subsection{Results}

Table~\ref{tab:empirical} reports the seven estimators under two trimming rules, separately for each binary threshold.

\begin{table}[h!]
    \centering
    \resizebox{\textwidth}{!}{\input{../output/tables/empirical_duflo.tex}}
    \caption[Returns to schooling from INPRES school construction]{Returns to schooling from INPRES school construction. Sample: both sexes with non-missing log hourly wage, cohorts 1957--1962 and 1968--1972, $G^* \in \{0, 1\}$, $N = 22{,}414$ ($201$ districts). Panel A uses $D = \mathbf{1}\{\text{yeduc} \geq 6\}$ (primary completion); Panel B uses $D = \mathbf{1}\{\text{yeduc} \geq 12\}$ (high-school completion). Outcome $Y$ is log hourly wage. Dictionary of $p = 45$ regressors from sex, birth-year dummies, Java, urban, and pairwise interactions (Appendix~\ref{sec:6c}). DML estimators use Lasso with $L = 5$-fold observation-level cross-fitting and the FSIF correction of Proposition~\ref{prop:2}. Superscript $s$ denotes symmetric trimming $\hat e(X) \in [0.10, 0.90]$ (Crump et al.\ 2009); superscript $\hat\alpha$ denotes the paper's one-sided data-driven rule of Conjecture~\ref{cor:onesided}, which returns $\hat\alpha = 0$ on this dictionary. Standard errors clustered at the district of birth. Bold marks the preferred estimate in each panel. Per-year CI rescales the per-complier interval by the threshold in years. Four-learner robustness in Appendix~\ref{sec:app_empirical}.}\label{tab:empirical}
\end{table}

\FloatBarrier
Two results organize the table. First, the FSIF correction systematically debiases the Lasso plug-in. Across the four plug-in–DML pairs (Wald$_X$ against DML-Wald and TC$_X$ against DML-TC, at each threshold), the plug-in overstates the return to schooling by between $0.2$ and $0.8$ standard errors, and the sign is always positive. At primary completion TC$_X = 1.27$ drops to DML-TC $= 1.06$ and Wald$_X = 2.20$ drops to DML-Wald $= 2.07$; the high-school corrections are smaller but point the same way. Cluster-robust intervals tighten by $6$ to $8\%$ at primary completion and are essentially unchanged at high school. The heterogeneity story between the two estimands is secondary. The DML-Wald and DML-TC gap of $2.07$ versus $1.06$ at primary completion cannot be blamed on a weak instrument (first-stage $F = 59.8$, and $48.2$ at high school) and is consistent with heterogeneous returns across districts with different INPRES exposure. DML-TC is preferred because it replaces the stronger stable-conditional-treatment-effects restriction in Assumption~\ref{ass:5} with the within-treatment-status common-trends restriction in Assumption~\ref{ass:4prime}.

Second, the trimming rule is inactive here. The $\hat\alpha$-rule returns $\hat\alpha = 0$, and the DML point estimates remain stable as $\alpha$ is pushed to $0.20$ (Appendix Figure~\ref{fig:sensitivity_app}). The propensity distribution is interior with support $[0.53, 0.91]$ and no mass near the boundary (Appendix Figure~\ref{fig:propensity_app}), which is what delivers the insensitivity. The $p = 45$ dictionary is deliberately built from individual-level regressors only, with no district aggregates that would mechanically pin down $G^*$.\footnote{In a companion exercise using a $p = 141$ dictionary that adds district-level covariates, $\hat e(X)$ saturates at the upper boundary, the $\hat\alpha$-rule binds at $\hat\alpha > 0$, and the DML-TC point estimate remains within one standard error of the $p = 45$ result. Replication code is available at \url{https://github.com/andres-mena/DID_LATE}.}
